\documentclass[11pt]{article}

\usepackage[T1]{fontenc}
\usepackage[utf8]{inputenc}
\PassOptionsToPackage{breaklinks=true}{hyperref}
\usepackage{lmodern}
\usepackage[margin=1in]{geometry}
\usepackage{amsmath}
\usepackage{amssymb}
\usepackage{amsthm}
\usepackage{mathtools}
\usepackage{graphicx}
\usepackage{tikz}
\usetikzlibrary{arrows.meta,positioning,calc,backgrounds,decorations.pathmorphing}
\newcommand{\orcidicon}{\includegraphics[width=0.32cm]{logo-orcid.pdf}}
\newcommand{\orcidA}{\href{https://orcid.org/\orcidauthorA}{\orcidicon}}
\usepackage[dvipsnames]{xcolor}
\usepackage{booktabs}
\usepackage{array}
\usepackage{tabularx}
\usepackage[font=small,labelfont=bf]{caption}
\usepackage{microtype}
\usepackage{fancyhdr}
\usepackage{hyperref}
\usepackage{cleveref}
\usepackage{xurl}
\usepackage{enumitem}

\numberwithin{equation}{section}

\hypersetup{
  colorlinks=true,
  linkcolor={blue!60!black},
  citecolor={blue!60!black},
  urlcolor={blue!60!black}
}

\raggedbottom

\newcolumntype{Y}{>{\raggedright\arraybackslash}X}
\newcolumntype{L}[1]{>{\raggedright\arraybackslash}p{#1}}
\newcommand{\code}[1]{\nolinkurl{#1}}

\setlength{\emergencystretch}{3em}
\widowpenalty=10000
\clubpenalty=10000
\displaywidowpenalty=10000

\fancypagestyle{firstpage}{%
  \fancyhf{}
  \fancyfoot[C]{\fontsize{7.5}{8}\selectfont Automorph Inc., Wilmington, DE, USA \;\texttt{ioannis@automorph.io}\; \mbox{Zenodo DOI: 10.5281/zenodo.23088991}\; \textcopyright\ 2026\; \mbox{CC-BY 4.0}}
  \renewcommand{\headrulewidth}{0pt}
  \renewcommand{\footrulewidth}{0pt}
}

\usepackage{aliascnt}
\newtheorem{theorem}{Theorem}[section]
\crefname{theorem}{Theorem}{Theorems}
\newaliascnt{proposition}{theorem}
\newtheorem{proposition}[proposition]{Proposition}
\aliascntresetthe{proposition}
\crefname{proposition}{Proposition}{Propositions}
\newaliascnt{lemma}{theorem}
\newtheorem{lemma}[lemma]{Lemma}
\aliascntresetthe{lemma}
\crefname{lemma}{Lemma}{Lemmas}
\newaliascnt{corollary}{theorem}
\newtheorem{corollary}[corollary]{Corollary}
\aliascntresetthe{corollary}
\crefname{corollary}{Corollary}{Corollarys}
\theoremstyle{definition}
\newaliascnt{definition}{theorem}
\newtheorem{definition}[definition]{Definition}
\aliascntresetthe{definition}
\crefname{definition}{Definition}{Definitions}
\newaliascnt{example}{theorem}
\newtheorem{example}[example]{Example}
\aliascntresetthe{example}
\crefname{example}{Example}{Examples}
\newaliascnt{observation}{theorem}
\newtheorem{observation}[observation]{Observation}
\aliascntresetthe{observation}
\crefname{observation}{Observation}{Observations}
\theoremstyle{remark}
\newaliascnt{remark}{theorem}
\newtheorem{remark}[remark]{Remark}
\aliascntresetthe{remark}
\crefname{remark}{Remark}{Remarks}

% Reduction-diagram styles
\tikzset{
  st/.style={inner sep=1.2pt, font=\small, fill=white, rounded corners=2pt},
  rw/.style={-{Stealth[length=5pt]}, thick},
  rws/.style={-{Stealth[length=5pt]}, thick, densely dashed},
  cellfill/.style={fill=blue!10, draw=none},
  defect/.style={draw=red!70!black, very thick, densely dotted},
  nf/.style={inner sep=1.5pt, font=\small, draw, fill=white, rounded corners=2pt},
}
\newcommand{\rto}{\to}
\newcommand{\rtos}{\mathrel{\to^{*}}}
\newcommand{\joinable}{\mathrel{\downarrow}}


\title{To Flatten a Stone with Six Birds:\\[4pt]
Critical Pairs, Holonomy, and Confluence\\
in Rewriting Systems}
\newcommand{\orcidauthorA}{0009-0009-7659-5964}
\author{Ioannis Tsiokos\,\orcidA{}}
\date{v1: 17 March 2026\qquad v2: 1 October 2026}

\begin{document}

\maketitle
\thispagestyle{firstpage}

\begin{abstract}
\noindent
Confluence says that the order in which rewrite steps are applied does not matter in the end; flatness says that, locally, transport does not depend on the route taken. This paper makes the analogy precise at the local level and tests it computationally. The natural carrier is not the cycle structure of the reduction graph, and we show by small counterexamples that cycle-based notions of graph flatness neither imply nor are implied by confluence. Instead, we record each branching of two one-step rewrites out of a common state, the configuration tested by local confluence. A branching has trivial \emph{elementary holonomy} when its two branches can be completed to a common descendant, so that its boundary closes up, and nontrivial holonomy when they cannot. With this definition, elementary flatness is exactly local confluence. Under termination, Newman's lemma turns the local statement into a global one: confluence, elementary flatness, and uniqueness of normal forms coincide. These results hold for arbitrary abstract rewriting systems and are checked in Lean~4 without finiteness assumptions. In string and term rewriting, the branchings that can fail are the rule overlaps recorded by critical pairs; we work with the concrete critical-pair instances that occur in terms reachable from chosen start terms. An exhaustive enumeration of all $66{,}066$ rewriting relations on one to four labeled states agrees with the theorems and yields smallest counterexamples showing that each hypothesis is needed. Curated string and left-linear term-rewriting examples, including an overlap below the root, show reachable critical-pair defects coinciding with the observed obstructions. Finally, two before-and-after examples read a completion step as holonomy elimination: the added rule leaves the branching in place but supplies the missing filler.
\end{abstract}

\section{Introduction}
\label{sec:introduction}

Confluence is rewriting theory's cleanest order-independence principle: if a state reduces in two different ways, the divergence should be joinable \cite{BaaderNipkow1998,Huet1980,DershowitzJouannaud1990,Terese2003}. Discrete transport geometry packages the same intuition in a different language: transport should be independent of route, and the obstruction is holonomy \cite{Meneses2021,BaezMuniain1994}. This paper argues that their local cores coincide once rewriting is given its missing two-dimensional cells. The relevant object is not the bare reduction graph. It is a reduction 2-complex whose $0$-cells are states, whose $1$-cells are one-step rewrites, and whose generating $2$-cells are local peaks or, in structured rewriting settings, reachable critical pairs \cite{Mimram2023,GuiraudMalbos2018,GuiraudMalbos2012}.

The first temptation is also the first mistake. If one tries to read flatness directly from the bare reduction graph, one obtains surrogates that are vacuous or misleading in the terminating settings studied here. \Cref{sec:bare-graph-not-enough} shows this explicitly: directed-loop flatness collapses to acyclicity on the curated terminating examples, while ordinary undirected graph cycle structure is not diagnostic for confluence. The negative result is therefore constructive. It tells us where the geometric data must live: not on graph loops alone, but on elementary branching cells.

Once those missing $2$-cells are restored, the local theorem-level core becomes clean. \Cref{sec:core-theorem} proves that elementary flatness is equivalent to local confluence. The paper then shows computationally that the corrected local/global picture survives exhaustive finite-ARS stress tests in the terminating regime, curated bounded string-rewriting examples, and a curated left-linear term-rewriting bridge with genuine nonroot overlap. In those structured settings, the relevant local generators are reachable critical pairs, and the before/after completion examples show that defect elimination can be read as holonomy elimination.

This paper is also a mathematics-facing contribution to the Six Birds line of work. The general closure vocabulary comes from \cite{Tsiokos2026SixBirdsFoundations}, while the treatment of route mismatch as a disciplined mathematical defect was developed earlier in \cite{Tsiokos2026CountStone}. The present paper does not repeat that general background. Instead, it uses that vocabulary to isolate a specific claim about rewriting: the local obstructions to confluence are carried by the generating $2$-cells on which elementary holonomy is tested.

The paper is organized around three levels of claim. First, there is a proved core: elementary flatness is equivalent to local confluence, and confluence plus normalization yields a unique normal form from a chosen start state. Second, there is a strong computational and structural bridge: exhaustive finite-ARS audits, curated string and term-rewriting examples, and explicit reduction 2-complex artifacts all support the corrected local geometric picture. Third, there is a deferred theorem package: the full finite terminating left-linear term-rewriting formulation in the exact target form remains to be written as a general theorem rather than a curated bridge. \Cref{sec:proved-supported-deferred} makes this boundary explicit.

\paragraph{Contributions.}
This paper makes four contributions.
\begin{enumerate}
    \item It identifies the correct geometric object for rewriting: a reduction 2-complex rather than a bare reduction graph.
    \item It proves the local formal core that elementary flatness is exactly local confluence (\Cref{thm:elementary-flatness-iff-local-confluence}), together with the supporting uniqueness-of-normal-form fact under confluence and normalization (\Cref{prop:confluence-normalization-unique-normal-form}).
    \item It shows computationally that the corrected local/global picture survives exhaustive finite-ARS stress tests and curated string and left-linear term-rewriting examples, while graph-only flatness surrogates fail (\Cref{sec:bare-graph-not-enough,sec:computational-bridge}).
    \item It gives an explicit operational reading of completion as holonomy elimination in before/after string and term-rewriting pairs (\Cref{sec:completion-holonomy}).
\end{enumerate}

The rest of the paper proceeds as follows. \Cref{sec:preliminaries-scope} fixes the paper's scope and its minimal Six Birds placement. \Cref{sec:bare-graph-not-enough} rules out the wrong graph-only formulation. \Cref{sec:reduction-two-complex} introduces the reduction 2-complex and elementary holonomy language, and \Cref{sec:core-theorem} proves the theorem-level core. \Cref{sec:critical-pairs-curvature} identifies reachable critical pairs as the structured rewriting generators of local curvature, \Cref{sec:computational-bridge} records the finite-ARS, string, and term-rewriting evidence base, and \Cref{sec:completion-holonomy} packages the before/after examples as holonomy-elimination case studies. \Cref{sec:proved-supported-deferred} then states explicitly what is proved, what is strongly supported, and what remains final theorem-writing work.

\section{Preliminaries and Scope}
\label{sec:preliminaries-scope}

\subsection{Abstract rewriting}
\label{subsec:abstract-rewriting}

An \emph{abstract rewriting system} is a pair $(X,\to)$ consisting of a set $X$ of states and a binary relation ${\to} \subseteq X \times X$, the one-step rewrite relation. No finiteness is assumed. We write $\to^{*}$ for the reflexive--transitive closure of $\to$, and say that $b$ is \emph{reachable} from $a$, or is a \emph{descendant} of $a$, when $a \to^{*} b$. The following notions are standard \cite{BaaderNipkow1998,Huet1980,Terese2003}.

\begin{definition}[Basic notions]
\label{def:basic-notions}
Let $(X,\to)$ be an abstract rewriting system.
\begin{enumerate}
    \item States $b,c \in X$ are \emph{joinable}, written $b \joinable c$, if there is a $d \in X$ with $b \to^{*} d$ and $c \to^{*} d$.
    \item A state $n$ is a \emph{normal form} if there is no $m$ with $n \to m$.
    \item The system is \emph{locally confluent} if $a \to b$ and $a \to c$ imply $b \joinable c$.
    \item The system is \emph{confluent} if $a \to^{*} b$ and $a \to^{*} c$ imply $b \joinable c$.
    \item The system is \emph{normalizing} if every state reaches some normal form, and \emph{terminating} if there is no infinite rewrite sequence $a_0 \to a_1 \to a_2 \to \cdots$.
    \item A state $a$ has \emph{unique reachable normal forms} if any two normal forms reachable from $a$ are equal. This asserts that there is \emph{at most} one; existence is a separate matter.
\end{enumerate}
\end{definition}

Termination implies normalization, but not conversely; and neither confluence nor uniqueness of normal forms implies either of them. When $X$ is finite, the system is the same thing as a finite directed graph, possibly with self-loops, and termination amounts to the absence of directed cycles.

In string and term rewriting, $X$ is a set of words or first-order terms and $\to$ is generated by a set of rules applied at any position and, for terms, under any substitution. These settings have additional structure, notably critical pairs, which we recall in \Cref{sec:critical-pairs-curvature}.

\subsection{Place in the Six Birds program}
\label{subsec:six-birds-placement}

\begin{remark}[Six Birds primitives]
\label{rem:six-birds-primitives}
This paper uses the closure vocabulary introduced in \cite{Tsiokos2026SixBirdsFoundations}, whose six primitives are operator rewrite (P1), constraints or feasibility gating (P2), route mismatch or holonomy (P3), staging (P4), packaging (P5), and accounting or audit (P6). The present paper is primarily about P3: its subject is local route mismatch in rewriting systems. P5 and P6 play supporting roles, in packaging rewriting data into reduction 2-complexes and in keeping every defect explicit in machine-readable records. The mathematics below does not depend on this vocabulary.
\end{remark}

Within that program, the paper is a mathematical companion to \cite{Tsiokos2026CountStone}, which treats route mismatch as a quantitative defect, and it draws on the finite-state and coarse-graining background of \cite{Tsiokos2026LayStone}. Neither is needed to read what follows.

\subsection{Three kinds of evidence}
\label{subsec:scope-evidence}

The paper combines three kinds of evidence, and we keep them separate throughout.

\begin{enumerate}
    \item \emph{Proved results.} The statements in \Cref{sec:core-theorem} are theorems about arbitrary abstract rewriting systems, with paper proofs and machine-checked Lean proofs. \Cref{prop:bare-graph-not-diagnostic} is proved by explicit counterexamples.
    \item \emph{Cited classical results.} The critical-pair lemma of Knuth--Bendix and Huet, which connects critical pairs to local confluence in string and term rewriting, is used as stated in the literature. We neither re-prove nor mechanize it.
    \item \emph{Computations.} An exhaustive enumeration of small finite systems, a search for smallest counterexamples, and a set of curated string and term-rewriting examples. Statements resting on these computations concern only the systems examined. Where a computation would need an unbounded search, as with term rewriting from a start term, a verdict is reported only after the set of reachable terms has been explored completely.
\end{enumerate}

The paper does not claim a new confluence criterion for term rewriting, nor a holonomy theory in the sense of a group-valued connection. Its contribution is a precise local dictionary, the global theorems that follow from it, and a computational test of both.

\section{Why Graph Cycles Are the Wrong Test}
\label{sec:bare-graph-not-enough}

For a finite abstract rewriting system the reduction graph \emph{is} the system: its vertices are the states and its directed edges are the one-step rewrites. Every property in \Cref{def:basic-notions}, confluence included, is therefore a property of that directed graph, and an algorithm can decide it from the graph by computing reachability. The question of this section is narrower. Graph theory and discrete geometry offer ready-made notions of flatness built from \emph{cycles}, and one may hope that one of them captures confluence. We show that the two most natural ones do not.

\begin{enumerate}
    \item \emph{Directed-cycle flatness} calls the reduction graph flat when it has no directed cycle.
    \item \emph{Undirected cycle rank} measures the dimension of the cycle space of the underlying simple undirected graph, i.e.\ the number of independent undirected cycles; rank zero means the graph is a forest.
\end{enumerate}

\begin{figure}[t]
\centering
\begin{tikzpicture}
  \begin{scope}
    \node[st] (A) at (0,0) {$A$};
    \node[st] (B) at (-0.9,-1.0) {$B$};
    \node[st] (C) at (0.9,-1.0) {$C$};
    \node[st] (D) at (0,-2.0) {$D$};
    \draw[rw] (A)--(B); \draw[rw] (A)--(C); \draw[rw] (B)--(D); \draw[rw] (C)--(D);
    \node[font=\small, align=center] at (0,-2.9) {(a) diamond $E001$};
    \node[font=\scriptsize, align=center] at (0,-3.6) {acyclic, cycle rank $1$\\ confluent};
  \end{scope}
  \begin{scope}[xshift=4.6cm]
    \node[st] (A) at (0,0) {$A$};
    \node[st] (B) at (-0.9,-1.0) {$B$};
    \node[nf] (C) at (0.9,-1.0) {$C$};
    \node[nf] (D) at (-0.9,-2.0) {$D$};
    \draw[rw] (A)--(B); \draw[rw] (A)--(C); \draw[rw] (B)--(D);
    \draw[defect] (D) .. controls (0.3,-2.2) and (0.9,-1.8) .. (C);
    \node[font=\small, align=center] at (0,-2.9) {(b) fork $E002$};
    \node[font=\scriptsize, align=center] at (0,-3.6) {acyclic, cycle rank $0$\\ not confluent};
  \end{scope}
  \begin{scope}[xshift=9.2cm]
    \node[st, minimum size=14pt] (s) at (0,-1.2) {$s$};
    \draw[rw] (s.north east) .. controls (0.9,0.2) and (-0.9,0.2) .. (s.north west);
    \node[font=\small, align=center] at (0,-2.9) {(c) one self-loop};
    \node[font=\scriptsize, align=center] at (0,-3.6) {directed cycle\\ confluent};
  \end{scope}
\end{tikzpicture}
\caption{Cycle structure says nothing about confluence. (a) The diamond is confluent and has an undirected cycle. (b) The fork has no cycles of either kind but is not confluent: its two normal forms $C$ and $D$ (boxed) have no common descendant, as indicated by the dotted line. (c) A single state with a self-loop is confluent, since every state reachable from $s$ is $s$ itself, yet it has a directed cycle.}
\label{fig:bare-graph-failure}
\end{figure}

\begin{table}[t]
    \caption{Cycle-based proxies on the curated finite systems $E001$--$E004$. All four graphs are acyclic, so directed-cycle flatness calls all of them flat, yet $E002$ and $E004$ are not confluent. Undirected cycle rank is positive on the confluent systems $E001$ and $E003$ and zero on the nonconfluent ones. The column of nonjoinable peaks is the local quantity that does track confluence (\Cref{sec:core-theorem}).}
    \label{tab:bare-graph-flatness-mismatch}
    \centering
    \smallskip
    \small
\begin{tabular}{@{}lcccc@{}}
\toprule
Example & \shortstack{Directed-cycle\\flat?} & \shortstack{Undirected\\cycle rank} & \shortstack{Nonjoinable\\peaks} & Confluent? \\
\midrule
$E001$ (diamond)           & yes & 1 & 0 & yes \\
$E002$ (fork)              & yes & 0 & 1 & no \\
$E003$ (delayed join)      & yes & 1 & 0 & yes \\
$E004$ (double fork)       & yes & 0 & 2 & no \\
\bottomrule
\end{tabular}

\end{table}

\begin{proposition}[Cycle-based proxies do not detect confluence]
\label{prop:bare-graph-not-diagnostic}
Among finite abstract rewriting systems, neither directed-cycle flatness nor vanishing undirected cycle rank implies confluence, and neither is implied by confluence. Directed-cycle flatness fails to imply confluence even among terminating systems, where it holds by definition.
\end{proposition}

\begin{proof}
All four claims are witnessed by the systems in \Cref{fig:bare-graph-failure}. The fork (b), with edges $A \to B$, $A \to C$, $B \to D$, has no directed cycle and its undirected graph is a tree, yet $C$ and $D$ are distinct normal forms reachable from $A$, so it is not confluent. The self-loop (c) has a directed cycle and is confluent. The diamond (a) has undirected cycle rank $1$ and is confluent, so confluence does not force rank zero. Finally, every terminating finite system is acyclic and hence directed-cycle flat, while the fork is terminating and not confluent.
\end{proof}

The exhaustive search of \Cref{subsec:boundary-search} finds these witnesses automatically: the smallest system on which directed-cycle flatness holds without confluence is the three-state fork $s_0 \to s_1$, $s_0 \to s_2$, and the smallest confluent system with a directed cycle is the one-state self-loop.

Why do cycle-based notions miss confluence? Because confluence compares pairs of rewrite sequences out of a common state, and locally, for instance in terminating systems where local confluence suffices, it is tested on \emph{peaks}, pairs of one-step rewrites $b \leftarrow a \to c$. A divergence need not contain a cycle, and the existence of a cycle does not determine whether its branches are joinable. In the fork the peak at $A$ cannot be joined, but nothing in the graph closes up around it; in the diamond the peak can be joined, and only then does it become part of a closed walk $A \to B \to D \leftarrow C \leftarrow A$. Cycles that do exist in a reduction graph may have nothing to do with peaks, as in the self-loop. A notion of flatness adapted to confluence must therefore first say \emph{which pairs of paths are being compared}. The reduction graph contains all the information needed to answer the question, but it does not single out the peaks as the objects on which the comparison takes place. The next section makes that comparison data explicit.

\section{Reduction 2-Complexes and Elementary Holonomy}
\label{sec:reduction-two-complex}

This section introduces the object on which the flatness reading of confluence is literal. It adds to the reduction graph a list of \emph{branchings}, the pairs of one-step rewrites to be compared, and fills a branching with a $2$-cell exactly when its two branches can be brought together.

\subsection{Branchings and fillers}

\begin{definition}[Elementary peaks and fillers]
\label{def:elementary-peak-flatness}
\label{def:boundary-filler-data}
Let $(X,\to)$ be an abstract rewriting system.
\begin{enumerate}
    \item An \emph{elementary peak} is a triple $(a,b,c)$ with $a \to b$, $a \to c$, and $b \neq c$. We call $a$ its \emph{source} and $b$, $c$ its \emph{branches}. The triples $(a,b,c)$ and $(a,c,b)$ describe the same \emph{branching}; when branchings are counted or cells attached, the two are identified and one representative is used.
    \item A \emph{filler} for the peak $(a,b,c)$ is a state $w$, the \emph{join witness}, together with rewrite sequences $b \to^{*} w$ and $c \to^{*} w$.
\end{enumerate}
\end{definition}

A filler turns the open shape $b \leftarrow a \to c$ into a closed walk
\[
a \to b \to^{*} w \leftarrow^{*} c \leftarrow a,
\]
which can serve as the boundary of a disk. Without a filler there is no closed boundary to attach anything to. This asymmetry is the whole point of the construction, and we build it into the definition.

\begin{definition}[Reduction 2-complex]
\label{def:reduction-two-complex}
Let $(X,\to)$ be an abstract rewriting system and let $G$ be a set of branchings, i.e.\ of elementary peaks up to exchange of the two branches, called the \emph{generating branchings}. The \emph{reduction 2-complex} of $(X,\to)$ with respect to $G$ consists of
\begin{enumerate}
    \item the states of $X$ as $0$-cells;
    \item the one-step rewrites $a \to b$ as directed $1$-cells;
    \item for each branching in $G$ that admits a filler, one $2$-cell attached along the closed walk of a chosen filler.
\end{enumerate}
Branchings in $G$ that admit no filler are kept as \emph{open branchings}: they are recorded, but no cell is attached to them.
\end{definition}

For an abstract rewriting system, $G$ is the set of all branchings. For string and term rewriting, \Cref{sec:critical-pairs-curvature} takes $G$ to be the set of critical-pair instances, the branchings that can actually fail. Whether a branching admits a filler does not depend on which filler is chosen; the choice affects only which $2$-cell is drawn. The computations of \Cref{sec:computational-bridge} fix a canonical filler by a deterministic rule (\Cref{app:computational-protocol}) so that the recorded complexes are reproducible.

\subsection{Elementary holonomy}

\begin{definition}[Elementary holonomy and elementary flatness]
\label{def:elementary-holonomy-status}
An elementary peak has \emph{trivial elementary holonomy} if it admits a filler, and \emph{nontrivial elementary holonomy} otherwise; the status is the same for $(a,b,c)$ and $(a,c,b)$, so it is a property of the branching. A rewriting system is \emph{elementarily flat} if every elementary peak has trivial elementary holonomy, that is, if the branches of every elementary peak are joinable. A reduction 2-complex records this status for the branchings in $G$: the filled ones carry $2$-cells and the others remain open.
\end{definition}

The term is chosen to evoke, not to import, the geometric notion. Elementary holonomy is a yes/no status attached to a single branching: either going down one branch and then on to $w$ agrees with going down the other branch and then on to $w$, for some $w$, or no such $w$ exists. It is not the holonomy of a connection valued in a group, and nothing below is a statement about fundamental groups, homotopy of rewrite paths, or higher coherence. Those richer questions are the subject of the homotopical theory of rewriting initiated by Squier. For a convergent (terminating and confluent) presentation of a monoid, Squier showed that a choice of confluence diagrams for the critical branchings yields a homotopy basis of the associated space of derivations, and the polygraphic setting extends this to higher dimensions \cite{Squier1987,SquierOtto1987,GuiraudMalbos2012,GuiraudMalbos2018,Mimram2023}. Elementary holonomy records only the first, Boolean layer of that structure: whether a given $2$-cell exists at all.

\begin{remark}[Comparing normal forms is a different test]
\label{rem:holonomy-vs-formal-core}
For terminating systems it is tempting to declare a peak defective when its two branches have different sets of reachable normal forms. This is not the same as nonjoinability, and it is not local. In the system of \Cref{fig:nf-mismatch}, the peak at $A$ is joinable at $D$, so its holonomy is trivial, yet $B$ reaches the normal forms $\{D,E\}$ while $C$ reaches only $\{D\}$. The discrepancy is caused by a different peak, the one at $B$, which is genuinely nonjoinable. In our computations, normal-form comparison is therefore recorded only as a secondary diagnostic, and only for terminating systems; elementary holonomy and all defect counts use joinability alone.
\end{remark}

\begin{figure}[t]
\centering
\begin{tikzpicture}
  \node[st] (A) at (0,0) {$A$};
  \node[st] (B) at (-1.1,-1.1) {$B$};
  \node[st] (C) at (1.1,-1.1) {$C$};
  \node[nf] (D) at (0,-2.2) {$D$};
  \node[nf] (E) at (-2.2,-2.2) {$E$};
  \begin{scope}[on background layer]
    \fill[cellfill] (A.center) -- (B.center) -- (D.center) -- (C.center) -- cycle;
  \end{scope}
  \draw[rw] (A)--(B); \draw[rw] (A)--(C);
  \draw[rw] (B)--(D); \draw[rw] (C)--(D); \draw[rw] (B)--(E);
  \draw[defect] (E) .. controls (-1.4,-2.8) and (-0.6,-2.8) .. (D);
  \node[font=\scriptsize, align=left, anchor=west] at (2.0,-0.6) {peak at $A$: joinable at $D$ (trivial)};
  \node[font=\scriptsize, align=left, anchor=west] at (2.0,-1.1) {normal forms: $B \mapsto \{D,E\}$, $C \mapsto \{D\}$};
  \node[font=\scriptsize, align=left, anchor=west] at (2.0,-1.6) {peak at $B$: $D$, $E$ not joinable (nontrivial)};
\end{tikzpicture}
\caption{Joinability and normal-form agreement differ. The peak at $A$ is filled (shaded) because $B$ and $C$ both reach $D$, even though the sets of normal forms reachable from $B$ and from $C$ differ. The obstruction to confluence sits at the peak $D \leftarrow B \to E$.}
\label{fig:nf-mismatch}
\end{figure}

\subsection{Representative examples}

The same local question arises in each of the settings studied later: given two one-step rewrites from a common state, is there a filler? \Cref{tab:elementary-holonomy-summary} lists the generating branchings and their holonomy in six representative examples. In the finite systems $E001$ and $E002$, the generating branchings are the peaks of \Cref{fig:bare-graph-failure}. In the string pair $E007$--$E008$ and the term pair $TRS003$--$TRS004$ they are critical-pair instances; these pairs are discussed in \Cref{sec:critical-pairs-curvature,sec:completion-holonomy}.

\begin{table}[t]
    \caption{Generating branchings in six representative examples. A filled branching carries a $2$-cell attached along the closed walk through its join witness; an open branching has nontrivial elementary holonomy and no attached cell.}
    \label{tab:elementary-holonomy-summary}
    \centering
    \smallskip
    \small
\begin{tabular}{@{}llcccp{0.20\linewidth}@{}}
\toprule
Example & Domain & 2-cells & Trivial & Nontrivial & Witness / defect \\
\midrule
$E001$ (diamond)       & finite ARS & 1 & 1 & 0 & D \\
$E002$ (fork)          & finite ARS & 1 & 0 & 1 & nonjoinable \\
$E007$ (before)        & string     & 1 & 0 & 1 & nonjoinable \\
$E008$ (after)         & string     & 1 & 1 & 0 & \texttt{y} \\
$TRS003$ (nested)      & term rew.  & 1 & 0 & 1 & nonjoinable \\
$TRS004$ (resolved)    & term rew.  & 1 & 1 & 0 & \texttt{p(a)} \\
\bottomrule
\end{tabular}

\end{table}

\section{Elementary Flatness, Local Confluence, and Termination}
\label{sec:core-theorem}

Throughout this section $(X,\to)$ is an arbitrary abstract rewriting system. No finiteness is assumed, and every result has a Lean counterpart listed in \Cref{app:lean-formal-core}.

\subsection{The local equivalence}

\begin{lemma}[Reflexive joinability]
\label{lem:joinable-refl}
Every state is joinable with itself.
\end{lemma}

\begin{proof}
Take $d = x$ and use $x \to^{*} x$ twice.
\end{proof}

\begin{theorem}[Elementary flatness is local confluence]
\label{thm:elementary-flatness-iff-local-confluence}
An abstract rewriting system is elementarily flat if and only if it is locally confluent.
\end{theorem}

\begin{proof}
Suppose the system is locally confluent and let $(a,b,c)$ be an elementary peak. Then $a \to b$ and $a \to c$, so $b \joinable c$ by local confluence. Hence every elementary peak has a filler.

Conversely, suppose the system is elementarily flat and let $a \to b$ and $a \to c$. If $b = c$, then $b \joinable c$ by \Cref{lem:joinable-refl}. If $b \neq c$, then $(a,b,c)$ is an elementary peak, which has a filler by assumption, so again $b \joinable c$.
\end{proof}

\begin{remark}[What the theorem says]
\label{rem:core-theorem-scope}
Elementary flatness is local confluence restricted to peaks with distinct branches, so the theorem is close to a restatement; its only extra content is the trivial case $b = c$. We state it as a theorem because it is the exact form of the dictionary between the two vocabularies: ``every elementary branching has trivial holonomy'' and ``every one-step divergence can be rejoined'' are the same condition. The theorem says nothing about which peaks need to be checked in a structured formalism; that is the role of critical pairs (\Cref{sec:critical-pairs-curvature}). Nor is it, by itself, a local-to-global statement. Termination supplies the global step.
\end{remark}

\subsection{From local to global under termination}

The remaining results relate the local condition to confluence and to normal forms. The first is standard.

\begin{proposition}[Confluence and normalization give a unique normal form]
\label{prop:confluence-normalization-unique-normal-form}
If $(X,\to)$ is confluent and normalizing, then every $a \in X$ reaches exactly one normal form.
\end{proposition}

\begin{proof}
Normalization gives a normal form $n$ with $a \to^{*} n$. If also $a \to^{*} m$ with $m$ a normal form, confluence gives a $d$ with $n \to^{*} d$ and $m \to^{*} d$. A normal form has no descendants other than itself, so $d = n$ and $d = m$; hence $n = m$.
\end{proof}

Uniqueness alone, the ``at most one'' half of this argument, uses only confluence. The next lemma is the familiar fact that termination yields normal forms. We use the formulation of termination as well-foundedness: every nonempty set of states contains an element with no one-step successor inside the set. This permits induction along $\to$, in which a property may be assumed for all one-step successors of $a$ when proving it for $a$.

\begin{lemma}[Termination implies normalization]
\label{lem:terminating-normalizing}
A terminating abstract rewriting system is normalizing.
\end{lemma}

\begin{proof}
By well-founded induction along $\to$. If $a$ is a normal form, it reaches itself. Otherwise $a \to b$ for some $b$; by the induction hypothesis $b$ reaches a normal form $n$, and so $a \to b \to^{*} n$.
\end{proof}

\begin{theorem}[Newman's lemma]
\label{thm:newman}
A terminating, locally confluent abstract rewriting system is confluent.
\end{theorem}

\begin{proof}
We prove by well-founded induction along $\to$ that every $a$ satisfies: whenever $a \to^{*} b$ and $a \to^{*} c$, we have $b \joinable c$. If either sequence is empty, say $a = b$, then $c$ itself is a common descendant of $b$ and $c$. Otherwise both sequences begin with a step, $a \to b_1 \to^{*} b$ and $a \to c_1 \to^{*} c$. By local confluence there is a $d$ with $b_1 \to^{*} d$ and $c_1 \to^{*} d$. Since $a \to c_1$, the induction hypothesis applies at $c_1$; it joins $c$ and $d$ at some $v$. Since $a \to b_1$ and $b_1 \to^{*} d \to^{*} v$, the induction hypothesis applies at $b_1$; it joins $b$ and $v$ at some $w$. Then $b \to^{*} w$ and $c \to^{*} v \to^{*} w$.
\end{proof}

\begin{figure}[t]
\centering
\begin{tikzpicture}[x=1.0cm,y=1.0cm]
  \node[st] (a) at (0,0) {$a$};
  \node[st] (b1) at (-1.2,-1.0) {$b_1$};
  \node[st] (c1) at (1.2,-1.0) {$c_1$};
  \node[st] (d) at (0,-2.0) {$d$};
  \node[st] (b) at (-2.6,-2.2) {$b$};
  \node[st] (c) at (2.6,-2.2) {$c$};
  \node[st] (v) at (1.3,-3.2) {$v$};
  \node[st] (w) at (-0.4,-4.3) {$w$};
  \begin{scope}[on background layer]
    \fill[cellfill] (a.center)--(b1.center)--(d.center)--(c1.center)--cycle;
    \fill[blue!5] (c1.center)--(d.center)--(v.center)--(c.center)--cycle;
    \fill[blue!16] (b1.center)--(b.center)--(w.center)--(v.center)--(d.center)--cycle;
  \end{scope}
  \draw[rw] (a)--(b1); \draw[rw] (a)--(c1);
  \draw[rws] (b1)--(b); \draw[rws] (c1)--(c);
  \draw[rws] (b1)--(d); \draw[rws] (c1)--(d);
  \draw[rws] (d)--(v); \draw[rws] (c)--(v);
  \draw[rws] (b)--(w); \draw[rws] (v)--(w);
  \node[font=\scriptsize] at (0,-1.0) {LC};
  \node[font=\scriptsize] at (1.3,-2.2) {IH at $c_1$};
  \node[font=\scriptsize] at (-1.0,-2.9) {IH at $b_1$};
\end{tikzpicture}
\caption{The tiling in the proof of Newman's lemma. Local confluence (LC) fills the elementary peak at $a$; the induction hypothesis (IH), available because $b_1$ and $c_1$ are one step below $a$, fills the two larger regions. Termination is what makes the induction well founded.}
\label{fig:newman-tiling}
\end{figure}

\Cref{fig:newman-tiling} displays the proof as a tiling: one elementary $2$-cell and two cells supplied by induction. This is the precise sense in which, for terminating systems, global flatness is assembled from elementary flatness. Combining \Cref{thm:newman} with \Cref{thm:elementary-flatness-iff-local-confluence}, and noting that confluence trivially implies local confluence, we obtain the global form of the dictionary.

\begin{corollary}[Confluence is elementary flatness under termination]
\label{cor:terminating-confluence-iff-flat}
A terminating abstract rewriting system is confluent if and only if it is elementarily flat.
\end{corollary}

The relation between confluence and normal forms needs less than termination.

\begin{theorem}[Confluence and unique normal forms]
\label{thm:normalizing-confluence-iff-unf}
Let $(X,\to)$ be normalizing. Then $(X,\to)$ is confluent if and only if every state has unique reachable normal forms.
\end{theorem}

\begin{proof}
If the system is confluent, two normal forms reachable from the same state are joinable and hence equal, as in the proof of \Cref{prop:confluence-normalization-unique-normal-form}. Conversely, suppose every state has unique reachable normal forms, and let $a \to^{*} b$ and $a \to^{*} c$. By normalization, $b \to^{*} n$ and $c \to^{*} m$ for normal forms $n$ and $m$. Both are reachable from $a$, so $n = m$, and this common normal form joins $b$ and $c$.
\end{proof}

\begin{corollary}[Existence and uniqueness of normal forms]
\label{cor:terminating-flat-unique-nf}
In a terminating, elementarily flat abstract rewriting system, every state reaches exactly one normal form.
\end{corollary}

\begin{proof}
By \Cref{cor:terminating-confluence-iff-flat} the system is confluent, by \Cref{lem:terminating-normalizing} it is normalizing, and \Cref{prop:confluence-normalization-unique-normal-form} applies.
\end{proof}

For terminating systems, then, four conditions coincide: elementary flatness, local confluence, confluence, and uniqueness of reachable normal forms. Each hypothesis used along the way is needed. \Cref{subsec:boundary-search} exhibits smallest finite witnesses: a four-state system that is locally confluent but not confluent, for which termination fails (\Cref{fig:newman-needs-termination}); a one-state system that is confluent but not normalizing; and a three-state system in which every state has unique reachable normal forms but which is not confluent, for which normalization fails.

\begin{remark}[Mechanization]
\label{rem:lean-formal-core-crossref}
All results in this section are proved in Lean~4 for an arbitrary relation on an arbitrary type, with no finiteness assumption and with termination stated as well-foundedness of the reduction relation in the direction of rewriting. The proofs use no axioms beyond Lean's standard ones (propositional extensionality, quotient soundness, and choice), and several use none. \Cref{app:lean-formal-core} gives the correspondence.
\end{remark}

\section{Critical Pairs as Generators of Holonomy}
\label{sec:critical-pairs-curvature}

In string and term rewriting, the one-step relation is generated by finitely many rules applied anywhere, so a single word or term typically has many peaks. Elementary flatness asks for all of them to be joinable. This section recalls why only a small, structured part of them can fail, namely the overlaps recorded by critical pairs, and explains how we use critical pairs relative to chosen start terms.

\subsection{Three kinds of peaks}
\label{subsec:three-kinds-of-peaks}

Consider a term $t$ with two redexes, one for a rule $l_1 \to r_1$ at position $p$ and one for a rule $l_2 \to r_2$ at position $q$. There are three possibilities \cite{Huet1980,BaaderNipkow1998,Terese2003}, illustrated in \Cref{fig:three-peaks}.

\begin{enumerate}
    \item \emph{Disjoint redexes.} Neither position lies below the other. The two rewrites act on separate parts of $t$ and commute: each result can be rewritten by the other rule to the same term. For example, with rules $a \to b$ and $c \to d$, the word $ac$ has the peak $bc \leftarrow ac \to ad$, joined at $bd$. No two left-hand sides overlap here.
    \item \emph{Variable overlap.} One redex lies inside the part of the other that is matched by a variable. The peak is again joinable: on the side where the outer rule has fired, the inner rewrite is applied to every copy of that variable's instance on the right-hand side; on the other side, it is applied to any remaining copies on the left-hand side, after which the outer rule fires. For left-linear rules there are no remaining copies. For example, with rules $f(x) \to g(x)$ and $a \to b$, the term $f(a)$ has the peak $g(a) \leftarrow f(a) \to f(b)$, joined at $g(b)$.
    \item \emph{Critical overlap.} One redex lies at a non-variable position of the other's left-hand side, so the two left-hand sides genuinely share symbols. Joinability is not automatic. Every such peak is an instance, inside some context and under some substitution, of a \emph{critical pair}: a pair of terms obtained by unifying one left-hand side with a non-variable subterm of another, using a most general unifier. A rule may overlap with a renamed copy of itself at a proper subterm; the overlap of a rule with itself at the root is excluded because it produces two identical terms.
\end{enumerate}

\begin{figure}[t]
\centering
\begin{tikzpicture}
  \begin{scope}
    \node[st] (s) at (0,0) {$ac$};
    \node[st] (l) at (-0.9,-1.0) {$bc$};
    \node[st] (r) at (0.9,-1.0) {$ad$};
    \node[st] (w) at (0,-2.0) {$bd$};
    \begin{scope}[on background layer]\fill[cellfill] (s.center)--(l.center)--(w.center)--(r.center)--cycle;\end{scope}
    \draw[rw] (s)--(l); \draw[rw] (s)--(r); \draw[rw] (l)--(w); \draw[rw] (r)--(w);
    \node[font=\small] at (0,-2.75) {(a) disjoint};
    \node[font=\scriptsize] at (0,-3.25) {$a \to b$,\; $c \to d$};
  \end{scope}
  \begin{scope}[xshift=4.4cm]
    \node[st] (s) at (0,0) {$f(a)$};
    \node[st] (l) at (-0.9,-1.0) {$g(a)$};
    \node[st] (r) at (0.9,-1.0) {$f(b)$};
    \node[st] (w) at (0,-2.0) {$g(b)$};
    \begin{scope}[on background layer]\fill[cellfill] (s.center)--(l.center)--(w.center)--(r.center)--cycle;\end{scope}
    \draw[rw] (s)--(l); \draw[rw] (s)--(r); \draw[rw] (l)--(w); \draw[rw] (r)--(w);
    \node[font=\small] at (0,-2.75) {(b) variable overlap};
    \node[font=\scriptsize] at (0,-3.25) {$f(x) \to g(x)$,\; $a \to b$};
  \end{scope}
  \begin{scope}[xshift=8.8cm]
    \node[st] (s) at (0,0) {$aba$};
    \node[nf] (l) at (-0.9,-1.0) {$aa$};
    \node[st] (r) at (0.9,-1.0) {$ab$};
    \node[nf] (w) at (0.9,-2.0) {$a$};
    \draw[rw] (s)--(l); \draw[rw] (s)--(r); \draw[rw] (r)--(w);
    \draw[defect] (l) .. controls (-0.7,-1.9) and (0.1,-2.2) .. (w);
    \node[font=\small] at (0,-2.75) {(c) critical overlap};
    \node[font=\scriptsize] at (0,-3.25) {$ab \to a$,\; $ba \to b$};
  \end{scope}
\end{tikzpicture}
\caption{The three kinds of peaks. (a) Rewrites at disjoint positions commute. (b) A rewrite inside the instance of a variable is absorbed after the outer rule fires. Both are always joinable. (c) Overlapping left-hand sides: in $aba$, the redexes $ab$ and $ba$ share the middle letter. This is the string example $E006$; its branches lead to the distinct normal forms $aa$ and $a$ (boxed), so the branching stays open.}
\label{fig:three-peaks}
\label{fig:string-overlap-nonconfluent}
\end{figure}

String rewriting is the special case of term rewriting with unary function symbols, and the same classification applies to words. Two left-hand sides overlap critically when a suffix of one is a prefix of the other, or when one occurs inside the other.

The first two cases are always joinable, so a nonjoinable peak must come from a critical overlap. Conversely, if every critical pair is joinable, then so is every instance of it in any context, because rewriting is closed under contexts and substitutions. This is the content of the classical critical-pair lemma.

\begin{theorem}[Critical-pair lemma \cite{KnuthBendix1970,Huet1980}]
\label{thm:critical-pair-lemma}
A term rewriting system is locally confluent if and only if all of its critical pairs are joinable. The same holds for string rewriting systems.
\end{theorem}

We use this theorem as stated in the literature (see also \cite[Chapter~6]{BaaderNipkow1998}); it is neither re-proved nor mechanized here. Combined with \Cref{thm:elementary-flatness-iff-local-confluence} and \Cref{cor:terminating-confluence-iff-flat}, it gives the holonomy form of the Knuth--Bendix criterion.

\begin{corollary}
\label{cor:cp-flatness}
A term or string rewriting system is elementarily flat if and only if all of its critical pairs are joinable. If the system is terminating, it is confluent if and only if all of its critical pairs are joinable.
\end{corollary}

In the language of \Cref{sec:reduction-two-complex}, critical pairs generate the nontrivial holonomy obligations: they are the branchings whose fillers must be checked, while disjoint and variable-overlap peaks are filled automatically. They are not the whole set of peaks. Even a joinable critical pair is a genuine generating branching, whose $2$-cell happens to be filled.

\subsection{Critical pairs relative to start terms}
\label{subsec:string-critical-pairs}
\label{subsec:reachable-critical-pairs-vs-peaks}

Our computations study a rewriting system from designated start terms. Let $S$ be the set of starts and $X_S$ the set of terms reachable from $S$. Every descendant of a term in $X_S$ lies in $X_S$, so the rewriting relation restricts to $X_S$, and two terms of $X_S$ are joinable in the restricted system exactly when they are joinable in the full one.

\begin{definition}[Reachable critical-pair instance]
\label{def:reachable-cp-instance}
A \emph{reachable critical-pair instance} is a peak $u \leftarrow t \to v$ with $t \in X_S$ and $u \neq v$ in which the two redexes form a critical overlap. It is recorded by its source $t$, its unordered pair of branches $\{u,v\}$, and the rules and positions involved.
\end{definition}

Two features of this definition matter. First, the instance is concrete. The source of a critical pair as computed from the rules is a most general overlap, such as $ab$ for the rules $ab \to x$ and $ab \to y$, or $f(g(x))$ for $f(g(x)) \to p(x)$ and $g(y) \to q(y)$. The terms actually reached from the starts contain instances of it: the word $zab$, or the ground term $f(g(a))$. Asking whether the general source itself is reachable misses these. Second, the definition depends on $S$. A critical pair may have no instance in $X_S$, in which case it plays no role in the restricted system.

The classification of \Cref{subsec:three-kinds-of-peaks} applies peak by peak, and the joins of disjoint and variable-overlap peaks stay inside $X_S$. Hence the restricted system on $X_S$ is locally confluent if and only if every reachable critical-pair instance is joinable. This is the start-relative form of \Cref{thm:critical-pair-lemma} that the computations test. The equivalence holds whether or not $X_S$ is finite. Our computational certification of it, however, requires complete exploration of $X_S$: the absence of a join witness in an incomplete bounded search does not certify nonjoinability (\Cref{subsec:certification}).

\subsection{Curated examples}
\label{subsec:left-linear-trs-bridge}

\Cref{tab:system-vs-reachable-cp-vs-graph-peak} compares, for the curated string and term examples, the number of critical pairs computed from the rules, the number of reachable critical-pair instances, and the number of peaks in the explored reduction graph. Since disjoint and variable-overlap peaks are joinable, the general relation is only that every \emph{nonjoinable} peak is a critical-pair instance. In the curated examples something stronger happens: every peak of the explored graph is a critical-pair instance and conversely. This is a property of these examples, which happen to contain no disjoint or variable-overlap peaks, and not a general fact.

\begin{table}[t]
    \caption{Critical pairs computed from the rules (``system''), reachable critical-pair instances, peaks in the explored reduction graph, and nonjoinable reachable instances. In every example the explored graph is complete, and the reachable instances coincide with the peaks.}
    \label{tab:system-vs-reachable-cp-vs-graph-peak}
    \centering
    \smallskip
    \small
\setlength{\tabcolsep}{4pt}
\begin{tabularx}{\linewidth}{@{}L{0.22\linewidth} L{0.11\linewidth} *{4}{>{\centering\arraybackslash}p{0.08\linewidth}} >{\centering\arraybackslash}p{0.11\linewidth}@{}}
\toprule
Example & Domain & \shortstack{System\\CPs} & \shortstack{Reach.\\CPs} & \shortstack{Graph\\peaks} & \shortstack{Defective\\CPs} & Confluent? \\
\midrule
$E005$ (sort)          & string    & 1 & 1 & 1 & 0 & yes \\
$E006$ (overlap)       & string    & 2 & 1 & 1 & 1 & no \\
$E007$ (before)        & string    & 1 & 1 & 1 & 1 & no \\
$E008$ (after)         & string    & 1 & 1 & 1 & 0 & yes \\
$TRS002$ (root)        & term rew. & 1 & 1 & 1 & 1 & no \\
$TRS003$ (nested)      & term rew. & 1 & 1 & 1 & 1 & no \\
$TRS004$ (resolved)    & term rew. & 1 & 1 & 1 & 0 & yes \\
\bottomrule
\end{tabularx}

\end{table}

The string example $E006$ (\Cref{fig:three-peaks}c) has rules $ab \to a$ and $ba \to b$, and the rules produce two critical pairs, from the overlaps $aba$ and $bab$. From the start word $aba$ the reachable words are $aba$, $aa$, $ab$, and $a$. Only the first critical pair has an instance there, and it is the single peak of the explored graph. Its branches $aa$ and $ab \to a$ end in different normal forms, so it is nonjoinable and the system is not confluent from $aba$. The second critical pair is irrelevant from this start.

\begin{observation}[Curated string examples]
\label{prop:reachable-critical-pairs-capture-string-obstructions}
In each of the string examples $E005$--$E008$, explored completely from its start word, the reachable critical-pair instances are exactly the peaks of the explored reduction graph, the nonjoinable instances are exactly the nonjoinable peaks, and the system is confluent from its start word exactly when there are none.
\end{observation}

The supporting computations are in \Cref{tab:system-vs-reachable-cp-vs-graph-peak,tab:curated-string-summary}. Each example is small enough to check by hand. In $E005$, the sorting rules $ba \to ab$, $ca \to ac$, $cb \to bc$ give, from $cba$, a single overlap of $cb$ and $ba$; both branches $bca$ and $cab$ reach $abc$. In $E006$, as above, the single reachable instance is nonjoinable. The pair $E007$--$E008$ is treated in \Cref{sec:completion-holonomy}.

The term examples test the same picture with first-order terms. In $TRS002$, the rules $f(a,y) \to g(y)$ and $f(x,b) \to h(x)$ overlap at the root of $f(a,b)$, giving the branches $g(b)$ and $h(a)$, which are distinct normal forms. In $TRS003$, the rules
\[
f(g(x)) \to p(x), \qquad g(a) \to b
\]
overlap below the root: the redex $g(a)$ sits at the argument position of $f(g(x))$, a non-variable position. The reachable instance has source $f(g(a))$ and branches $p(a)$ and $f(b)$, again distinct normal forms. In $TRS004$ the rule $f(b) \to p(a)$ is added; the same instance remains, but now $f(b) \to p(a)$ fills it.

\begin{observation}[Curated left-linear term examples]
\label{prop:trs-critical-pairs-as-curvature-generators}
In each of the term examples $TRS002$--$TRS004$, explored completely from its start term, the reachable critical-pair instances are exactly the peaks of the explored reduction graph, and the system is confluent from its start term exactly when every instance is joinable. In $TRS003$ and $TRS004$ the overlap lies below the root.
\end{observation}

These observations are small. Their purpose is to confirm on concrete inputs that the implementation finds the right objects: instances in context rather than general overlap sources, overlaps below the root as well as at it, and a verdict tied to joinability.

\section{Computational Bridge and Stress Tests}
\label{sec:computational-bridge}

The preceding sections isolate the corrected geometric object and the theorem-level local core. This section records how that picture survives the implemented stress tests. The evidence comes in three layers:
\begin{enumerate}
    \item an exhaustive audit of all labeled finite abstract rewriting systems on up to four states in the terminating regime;
    \item curated bounded string-rewriting examples;
    \item a curated left-linear term-rewriting bridge with genuine nonroot overlap.
\end{enumerate}
The role of this section is not to replace theorem-writing. Its role is to show that once the object and the local generators are chosen correctly, the expected local/global relationships remain stable across every computational regime implemented in the repository.

\subsection{Exhaustive finite-ARS audit}
\label{subsec:exhaustive-finite-ars-audit}

The cleanest stress test is the exhaustive finite-ARS audit. The search space contains every labeled directed rewriting relation on $n$ states for $n \leq 4$, including self-loops. This yields
\[
2^1 + 2^4 + 2^9 + 2^{16} = 66066
\]
systems in total. Of these, $572$ are terminating.

\begin{proposition}[Exhaustive finite-ARS stress test]
\label{prop:exhaustive-finite-ars-stress-test}
Within the terminating finite-ARS search space on up to four states, the four target discrepancy classes tracked in the repository are empty.
\end{proposition}

\begin{proof}[Audit summary]
The search checks $2$, $16$, $512$, and $65536$ labeled systems at sizes $1$, $2$, $3$, and $4$ respectively, for a total of $66066$. The terminating counts are $1$, $3$, $25$, and $543$, for a total of $572$ terminating systems. On this terminating subset, the discrepancy counts
\[
\begin{aligned}
&\texttt{local\_flatness\_vs\_local\_confluence},\\
&\texttt{terminating\_local\_flatness\_not\_confluent},\\
&\texttt{terminating\_confluent\_not\_unique\_reachable\_normal\_forms},\\
&\texttt{terminating\_unique\_reachable\_normal\_forms\_not\_confluent}
\end{aligned}
\]
are all zero. Equivalently, on the terminating finite-ARS regime explored here, the local peak picture survives every system in the search space while the naive graph-only flatness surrogates do not.
\end{proof}

\begin{table}[t]
\caption{Exhaustive finite-ARS audit on all labeled systems up to four states. On the terminating subset, all four discrepancy classes vanish.}
\label{tab:exhaustive-finite-ars-summary}
\centering
\smallskip
\small
\setlength{\tabcolsep}{5pt}
\begin{tabularx}{0.96\linewidth}{L{0.14\linewidth} *{6}{>{\centering\arraybackslash}X}}
\toprule
Scope & \shortstack{Total\\systems} & \shortstack{Terminating} & LF$\neq$LC & \shortstack{LF \&\\$\neg$Conf} & \shortstack{Conf \&\\$\neg$UNF} & \shortstack{UNF \&\\$\neg$Conf} \\
\midrule
$n \leq 4$ & 66\,066 & 572 & 0 & 0 & 0 & 0 \\
\bottomrule
\end{tabularx}
\end{table}


The same finite-ARS infrastructure was also used outside the safe theorem hypotheses. Within the same small search bound, the boundary-counterexample search found explicit witnesses to the weakened claims obtained by dropping termination or replacing joinability-based local flatness with naive graph or normal-form surrogates. The positive finite-ARS results are therefore not an artifact of an impoverished search space; they are specific to the correct hypotheses.

\subsection{Curated string systems}
\label{subsec:curated-string-systems}

The string layer provides the first structured rewriting bridge. All four curated string examples explored completely under their chosen bounds, and all four exhibit the expected relation between global confluence, local peak defects, and reachable critical-pair defects.

\begin{table}[t]
\caption{Curated string examples. The table reports explored graph size, reachable critical-pair counts, reachable defect counts, normal forms, and confluence status for $E005$--$E008$.}
\label{tab:curated-string-summary}
\centering
\smallskip
\small
\setlength{\tabcolsep}{4pt}
\begin{tabularx}{\linewidth}{@{}L{0.10\linewidth} *{2}{>{\centering\arraybackslash}p{0.07\linewidth}} *{2}{>{\centering\arraybackslash}p{0.10\linewidth}} Y >{\centering\arraybackslash}p{0.11\linewidth}@{}}
\toprule
Example & States & Edges & \shortstack{Reach.\\CPs} & \shortstack{Defective\\CPs} & Normal forms & Confluent? \\
\midrule
$E005$ & 6 & 6 & 1 & 0 & \texttt{abc} & yes \\
$E006$ & 4 & 3 & 1 & 1 & \texttt{a}, \texttt{aa} & no \\
$E007$ & 3 & 2 & 1 & 1 & \texttt{x}, \texttt{y} & no \\
$E008$ & 3 & 3 & 1 & 0 & \texttt{y} & yes \\
\bottomrule
\end{tabularx}
\end{table}


The sorting example $E005$ explores completely to a six-state, six-edge reduction graph with one local peak, one reachable critical pair, zero local defects, zero reachable critical-pair defects, and unique normal form \texttt{abc}. The overlap example $E006$ explores completely to a four-state, three-edge reduction graph with one local peak, two system critical pairs, one reachable critical pair, and normal forms \texttt{a} and \texttt{aa}; its local peak defect counts and reachable critical-pair defect counts are both nonzero, and the explored system is nonconfluent. The completion-before example $E007$ explores completely to a three-state, two-edge reduction graph with normal forms \texttt{x} and \texttt{y}; it has one local peak, one reachable critical pair, nonzero local defects, nonzero reachable critical-pair defects, and nonconfluence. The completion-after example $E008$ explores completely to a three-state, three-edge reduction graph with unique normal form \texttt{y}; it has one local peak, one reachable critical pair, zero local defects, zero reachable critical-pair defects, and confluence.

Taken together, the string examples do not merely show that some nonconfluent systems have bad critical pairs. They show the tighter statement needed for the paper: on the curated examples, the reachable critical-pair defect layer and the observed graph-peak defect layer agree exactly, and both align with the global confluence behavior of the explored system.

\subsection{Curated left-linear term-rewriting bridge}
\label{subsec:curated-trs-bridge}

The left-linear term-rewriting prototype extends the same local/global pattern to a genuine first-order setting. All four curated examples explore completely under their chosen bounds.

\begin{table}[t]
\caption{Curated left-linear term-rewriting bridge. The table reports explored graph size, reachable critical-pair counts, reachable defect counts, normal forms, and confluence status for $TRS001$--$TRS004$.}
\label{tab:curated-trs-summary}
\centering
\smallskip
\small
\setlength{\tabcolsep}{4pt}
\begin{tabularx}{\linewidth}{@{}L{0.12\linewidth} *{2}{>{\centering\arraybackslash}p{0.07\linewidth}} *{2}{>{\centering\arraybackslash}p{0.10\linewidth}} Y >{\centering\arraybackslash}p{0.11\linewidth}@{}}
\toprule
Example & States & Edges & \shortstack{Reach.\\CPs} & \shortstack{Defective\\CPs} & Normal forms & Confluent? \\
\midrule
$TRS001$ & 3 & 2 & 0 & 0 & \texttt{s(s(z))} & yes \\
$TRS002$ & 3 & 2 & 1 & 1 & \texttt{g(b)}, \texttt{h(a)} & no \\
$TRS003$ & 3 & 2 & 1 & 1 & \texttt{f(b)}, \texttt{p(a)} & no \\
$TRS004$ & 3 & 3 & 1 & 0 & \texttt{p(a)} & yes \\
\bottomrule
\end{tabularx}
\end{table}


The Peano-addition smoke test $TRS001$ explores completely to a three-state, two-edge reduction graph with unique normal form \texttt{s(s(z))} and no critical pairs. The root-overlap example $TRS002$ explores completely to a three-state, two-edge reduction graph with one reachable critical pair, one graph peak, normal forms \texttt{g(b)} and \texttt{h(a)}, and a nontrivial reachable defect. The nested-overlap nonconfluent example $TRS003$ explores completely to a three-state, two-edge reduction graph with one reachable critical pair, one graph peak, normal forms \texttt{f(b)} and \texttt{p(a)}, and a nontrivial reachable defect. The nested-overlap resolved example $TRS004$ explores completely to a three-state, three-edge reduction graph with one reachable critical pair, one graph peak, unique normal form \texttt{p(a)}, and zero reachable defects.

The term-rewriting bridge therefore reproduces the same operational pattern seen in the string layer. The local generator persists, but its defect status changes. In the defective case the explored system is nonconfluent from the chosen start term; in the resolved case the explored system is confluent. Moreover, the reachability cross-check between critical-pair signatures and graph-peak signatures is exact on all curated TRS examples.

\begin{remark}[What the computational bridge does and does not establish]
\label{rem:computational-bridge-scope}
The computational evidence in this section is strong enough to fix the correct local geometric picture. It is not, by itself, the final theorem package for finite terminating left-linear term rewriting systems in full generality. What it does establish is that the corrected local/global story survives every implemented stress test: exhaustive finite ARS search in the terminating regime, curated string overlaps and completion pairs, and curated left-linear term-rewriting examples with genuine nonroot overlap.
\end{remark}

\section{Completion as Holonomy Elimination}
\label{sec:completion-holonomy}

Knuth--Bendix completion repairs a nonconfluent system by orienting nonjoinable critical pairs into new rules, and, in its general form, also simplifies or removes existing rules \cite{KnuthBendix1970,DershowitzJouannaud1990,BaaderNipkow1998}. We consider only the simplest kind of step, the addition of one rule. In the two rule-addition steps studied here, the overlap that produced the branching is still present after the new rule is added, so the branching remains; what changes is that it acquires a filler. We illustrate this with one string example and one term example, in each case adding a single rule by hand.

\begin{figure}[t]
\centering
\begin{tikzpicture}
  % string before
  \begin{scope}
    \node[st] (s) at (0,0) {$ab$};
    \node[nf] (x) at (-0.8,-1.1) {$x$};
    \node[nf] (y) at (0.8,-1.1) {$y$};
    \draw[rw] (s)--(x); \draw[rw] (s)--(y);
    \draw[defect] (x) .. controls (-0.3,-1.7) and (0.3,-1.7) .. (y);
    \node[font=\small] at (0,-2.3) {$E007$: before};
  \end{scope}
  \begin{scope}[xshift=2.9cm]
    \node[st] (s) at (0,0) {$ab$};
    \node[st] (x) at (-0.8,-1.1) {$x$};
    \node[nf] (y) at (0.8,-1.1) {$y$};
    \begin{scope}[on background layer]\fill[cellfill] (s.center)--(x.center)--(y.center)--cycle;\end{scope}
    \draw[rw] (s)--(x); \draw[rw] (s)--(y);
    \draw[rw, blue!60!black] (x)--node[below, font=\scriptsize]{new} (y);
    \node[font=\small] at (0,-2.3) {$E008$: after};
  \end{scope}
  % term before
  \begin{scope}[xshift=6.6cm]
    \node[st] (s) at (0,0) {$f(g(a))$};
    \node[nf] (l) at (-0.9,-1.1) {$f(b)$};
    \node[nf] (r) at (0.9,-1.1) {$p(a)$};
    \draw[rw] (s)--(l); \draw[rw] (s)--(r);
    \draw[defect] (l) .. controls (-0.3,-1.7) and (0.3,-1.7) .. (r);
    \node[font=\small] at (0,-2.3) {$TRS003$: before};
  \end{scope}
  \begin{scope}[xshift=9.9cm]
    \node[st] (s) at (0,0) {$f(g(a))$};
    \node[st] (l) at (-0.9,-1.1) {$f(b)$};
    \node[nf] (r) at (0.9,-1.1) {$p(a)$};
    \begin{scope}[on background layer]\fill[cellfill] (s.center)--(l.center)--(r.center)--cycle;\end{scope}
    \draw[rw] (s)--(l); \draw[rw] (s)--(r);
    \draw[rw, blue!60!black] (l)--node[below, font=\scriptsize]{new} (r);
    \node[font=\small] at (0,-2.3) {$TRS004$: after};
  \end{scope}
\end{tikzpicture}
\caption{Completion as holonomy elimination. Left: the string rules $ab \to x$ and $ab \to y$ produce an open branching with distinct normal forms $x$ and $y$; adding $x \to y$ fills it with join witness $y$. Right: the term rules $f(g(x)) \to p(x)$ and $g(a) \to b$ produce an open branching at $f(g(a))$; adding $f(b) \to p(a)$ fills it with join witness $p(a)$. In both cases the source and the two branches are unchanged.}
\label{fig:string-completion-before-after}
\label{fig:trs-completion-before-after}
\end{figure}

\subsection{The two examples}
\label{subsec:string-completion-holonomy}
\label{subsec:trs-completion-holonomy}

The string system $E007$ has the rules $ab \to x$ and $ab \to y$, whose left-hand sides coincide. From the start word $ab$ there is one branching, $x \leftarrow ab \to y$, and two normal forms, so the branching has nontrivial holonomy and the system is not confluent from $ab$. The system $E008$ adds the rule $x \to y$. The branching is still there, with the same source and branches, but now $x \to y$ provides a filler with join witness $y$. Its holonomy is trivial, and $y$ is the unique normal form.

The term system $TRS003$ has the rules $f(g(x)) \to p(x)$ and $g(a) \to b$, which overlap below the root. From the start term $f(g(a))$ there is one branching, $p(a) \leftarrow f(g(a)) \to f(b)$, whose branches are distinct normal forms. The system $TRS004$ adds the rule $f(b) \to p(a)$, which provides a filler with join witness $p(a)$.

\begin{table}[t]
    \caption{Before and after the added rule. The number of branchings does not change; the number of nonjoinable ones drops to zero, and confluence from the start term is restored.}
    \label{tab:holonomy-elimination-deltas}
    \centering
    \smallskip
    \small
\setlength{\tabcolsep}{5pt}
\begin{tabular}{@{}lllcccccc@{}}
\toprule
 & & & \multicolumn{2}{c}{Branchings} & \multicolumn{2}{c}{Nonjoinable} & \multicolumn{2}{c}{Confluent?} \\
\cmidrule(lr){4-5}\cmidrule(lr){6-7}\cmidrule(lr){8-9}
Domain & Before $\to$ after & Added rule & before & after & before & after & before & after \\
\midrule
string & $E007 \to E008$ & $x \to y$ & 1 & 1 & 1 & 0 & no & yes \\
term & $TRS003 \to TRS004$ & $f(b) \to p(a)$ & 1 & 1 & 1 & 0 & no & yes \\
\bottomrule
\end{tabular}

\end{table}

\begin{observation}[Completion as holonomy elimination]
\label{prop:completion-as-holonomy-elimination}
In the string pair $E007$--$E008$ and the term pair $TRS003$--$TRS004$, explored completely from their start terms, the added rule
\begin{enumerate}
    \item leaves the reachable critical-pair instance in place, with the same source and branches;
    \item changes its elementary holonomy from nontrivial to trivial;
    \item reduces the number of nonjoinable reachable instances from one to zero; and
    \item makes the system confluent from its start term.
\end{enumerate}
\end{observation}

In both pairs the ``after'' system consists of exactly the ``before'' rules plus the one added rule, with the same start term, so that the comparison is between a system and a one-rule extension of it. For the string pair this is checked automatically as part of the completion computation; for the term pair it can be read off the two rule sets.

\subsection{What is eliminated}
\label{subsec:what-is-eliminated}

The examples separate two quantities that are easily conflated: the number of generating branchings and the number of nontrivial ones. Completion reduces the second while leaving the first unchanged. In the reduction 2-complex, the effect of the new rule is to add a $1$-cell along which a $2$-cell can now be attached; the branching that was open becomes the boundary of a filled cell.

\begin{remark}[Scope]
\label{rem:completion-boundary}
The two examples are hand-built single steps, chosen so that one added rule suffices. They illustrate a reading of completion; they are not a geometric formulation of the completion procedure. In general, an added rule can create new critical pairs, completion may fail or fail to terminate, and orienting an equation requires a reduction ordering \cite{KnuthBendix1970,BaaderNipkow1998}. None of this is addressed here.
\end{remark}

\section{What Is Proved, What Is Computed, What Remains Open}
\label{sec:proved-supported-deferred}

\Cref{tab:claim-support-matrix} summarizes the status of each claim. This section comments on the three groups.

\begin{table}[t]
    \caption{Status of the main claims. ``Proved'' means proved in the paper and, where indicated, in Lean. ``Cited'' means a classical result used as stated in the literature. ``Computed'' means verified on the specific systems examined. ``Open'' means not established here.}
    \label{tab:claim-support-matrix}
    \centering
    \smallskip
    \small
\setlength{\tabcolsep}{5pt}
\renewcommand{\arraystretch}{1.15}
\begin{tabularx}{\linewidth}{@{}Y L{0.30\linewidth} l@{}}
\toprule
Claim & Evidence & Status \\
\midrule
Elementary flatness $\iff$ local confluence (\Cref{thm:elementary-flatness-iff-local-confluence}) & paper proof; Lean & proved \\
Confluence and normalization give a unique normal form (\Cref{prop:confluence-normalization-unique-normal-form}) & paper proof; Lean & proved \\
Newman's lemma; under termination, confluence $\iff$ elementary flatness (\Cref{thm:newman}, \Cref{cor:terminating-confluence-iff-flat}) & paper proof; Lean & proved \\
Under normalization, confluence $\iff$ unique reachable normal forms (\Cref{thm:normalizing-confluence-iff-unf}) & paper proof; Lean & proved \\
Cycle-based graph flatness does not detect confluence (\Cref{prop:bare-graph-not-diagnostic}) & explicit counterexamples & proved \\
Critical-pair lemma and its holonomy form (\Cref{thm:critical-pair-lemma}, \Cref{cor:cp-flatness}) & literature \cite{KnuthBendix1970,Huet1980} & cited \\
Theorems agree with all terminating systems on $1$--$4$ states; each hypothesis needed (\Cref{tab:exhaustive-finite-ars-summary,tab:boundary-counterexamples}) & exhaustive enumeration & computed \\
Reachable critical-pair instances match the observed obstructions in the curated string and term examples (\Cref{prop:reachable-critical-pairs-capture-string-obstructions,prop:trs-critical-pairs-as-curvature-generators}) & complete exploration of curated examples & computed \\
A single added rule fills the open branching in two curated pairs (\Cref{prop:completion-as-holonomy-elimination}) & complete exploration of curated examples & computed \\
Mechanized critical-pair lemma for left-linear term rewriting in 2-complex form & --- & open \\
Group-valued or connection-based holonomy theory of confluence & --- & open \\
\bottomrule
\end{tabularx}

\end{table}

\subsection{Proved}
\label{subsec:what-is-proved}

For abstract rewriting systems the picture is complete. Elementary flatness is local confluence (\Cref{thm:elementary-flatness-iff-local-confluence}). Under termination, Newman's lemma makes it equivalent to confluence (\Cref{cor:terminating-confluence-iff-flat}); under the weaker hypothesis of normalization, confluence is equivalent to uniqueness of reachable normal forms (\Cref{thm:normalizing-confluence-iff-unf}); and a terminating, elementarily flat system has exactly one normal form reachable from each state (\Cref{cor:terminating-flat-unique-nf}). These are classical facts of rewriting theory restated in the vocabulary of holonomy, and all of them are checked in Lean for arbitrary, possibly infinite, systems. The failure of cycle-based graph flatness (\Cref{prop:bare-graph-not-diagnostic}) is proved by explicit counterexamples.

\subsection{Cited and computed}
\label{subsec:what-is-supported}

For string and term rewriting, the link between critical pairs and local confluence is the classical critical-pair lemma, which we cite (\Cref{thm:critical-pair-lemma}). Together with the results above it gives \Cref{cor:cp-flatness}: a system is elementarily flat if and only if its critical pairs are joinable, and, when terminating, confluent under the same condition.

The computations add evidence of a different kind. On finite systems with one to four states they agree with the theorems on every case and produce smallest counterexamples for each dropped hypothesis. On the curated string and term examples they show that the start-relative analysis, based on concrete critical-pair instances in reachable terms, finds exactly the observed obstructions, including an overlap below the root, and that a single added rule can turn an open branching into a filled one. These are statements about the examples examined, and we do not extrapolate them.

\subsection{Open}
\label{subsec:what-remains}

Three directions remain open.

\begin{enumerate}
    \item \emph{A mechanized critical-pair lemma.} The Lean development covers abstract rewriting only. A formal proof of \Cref{thm:critical-pair-lemma} and its start-relative form for left-linear term rewriting, phrased in terms of reduction 2-complexes, requires the full case analysis of \Cref{subsec:three-kinds-of-peaks} together with stability of rewriting under contexts and substitutions. It is not attempted here.
    \item \emph{A genuinely geometric formulation.} Elementary holonomy is a Boolean status. A theory in which holonomy takes values in a group, or in which confluence becomes the flatness of an independently defined connection on a reduction complex, would need definitions that this paper does not supply. The homotopical approach to rewriting \cite{Squier1987,GuiraudMalbos2012,GuiraudMalbos2018,Mimram2023} is the natural setting for such a development.
    \item \emph{Completion in general.} \Cref{sec:completion-holonomy} reads two single completion steps as holonomy elimination. Whether the full completion procedure, with its orderings, new critical pairs, and possible failure, admits a useful description in these terms is open.
\end{enumerate}

\section{Conclusion}
\label{sec:conclusion}

The analogy between confluence and flatness is exact once one compares the right things. Cycles in the reduction graph are the wrong things: they neither detect nor are detected by confluence. Confluence compares pairs of rewrite sequences from a common state; local confluence compares the two branches of each peak, and under termination the two conditions are equivalent. A peak becomes a closed boundary, and so something that can be filled, only when its branches have a common descendant. Recording peaks as generating branchings, and filling them when possible, turns the reduction graph into a reduction 2-complex on which elementary holonomy is defined. Elementary flatness is then exactly local confluence, and under termination it is exactly confluence. The proof of Newman's lemma reads as a tiling of a large region by elementary cells, and these statements are checked in Lean for arbitrary rewriting systems.

In string and term rewriting, critical pairs mark the branchings that can fail; disjoint and variable-overlap peaks are filled automatically. Working relative to chosen start terms, with concrete critical-pair instances and with verdicts certified only for completely explored reachable sets, the computations find exactly the expected obstructions on the curated examples. In two completion examples, the added rule leaves the branching in place and supplies its filler.

The central message can be stated without exaggeration. Local order-independence of reduction, the joinability of every peak, is exactly elementary flatness, and for terminating systems it is equivalent to full order-independence, that is, to confluence. On the reduction 2-complex, the local generators of nontrivial holonomy are the nonjoinable peaks, and in structured rewriting they are instances of critical pairs.


\appendix
\crefalias{section}{appendix}
\crefalias{subsection}{subappendix}

\section{Lean Formalization}
\label{app:lean-formal-core}

The results of \Cref{sec:core-theorem} are formalized in Lean~4 (version 4.28.0) in the package \texttt{ConfluenceFlat}, which has no external dependencies. The development covers abstract rewriting only; the string and term-rewriting computations are not formalized.

\subsection{Setting}

The basic structure is a rewriting system on an arbitrary type, given by a single relation \texttt{step}. There is no finiteness assumption, and every definition and theorem below quantifies over such systems. A separate structure for finite systems extends it with a finite list of states containing every state that occurs in a step; it is provided for representing concrete examples and is not used by any theorem. Reachability is the inductively defined reflexive--transitive closure of \texttt{step}. Joinability, normal forms, local confluence, elementary peaks (with distinct branches), elementary flatness, confluence, normalization, and unique reachable normal forms are defined exactly as in \Cref{def:basic-notions,def:elementary-peak-flatness,def:elementary-holonomy-status}. Termination is well-foundedness of the relation ``$b$ is a one-step successor of $a$'', so that well-founded induction proceeds along rewrite steps. Uniqueness of reachable normal forms means ``at most one''; existence and uniqueness together is a separate predicate.

\subsection{Correspondence}

\Cref{tab:lean-correspondence} lists the Lean theorem corresponding to each result. All names are in the namespace \texttt{ConfluenceFlat.FiniteARS}, a name retained for continuity although the theorems are not restricted to finite systems.

\begin{table}[h]
\caption{Results of \Cref{sec:core-theorem} and their Lean counterparts.}
\label{tab:lean-correspondence}
\centering
\small
\begin{tabular}{@{}ll@{}}
\toprule
Result & Lean theorem \\
\midrule
\Cref{lem:joinable-refl} & \texttt{joinable\_refl} \\
\Cref{thm:elementary-flatness-iff-local-confluence} & \texttt{elementaryFlat\_iff\_localConfluent} \\
\Cref{prop:confluence-normalization-unique-normal-form} & \texttt{existsUniqueNormalFormFrom\_of\_confluent\_normalizing} \\
\Cref{lem:terminating-normalizing} & \texttt{normalizing\_of\_terminating} \\
\Cref{thm:newman} & \texttt{confluent\_of\_terminating\_localConfluent} \\
\Cref{cor:terminating-confluence-iff-flat} & \texttt{confluent\_iff\_elementaryFlat\_of\_terminating} \\
\Cref{thm:normalizing-confluence-iff-unf} & \texttt{confluent\_iff\_uniqueNormalForms\_of\_normalizing} \\
\Cref{cor:terminating-flat-unique-nf} & \texttt{existsUniqueNormalFormFrom\_of\_terminating\_elementaryFlat} \\
\bottomrule
\end{tabular}
\end{table}

Supporting lemmas include \texttt{reachable\_eq\_of\_normalForm} (a normal form reaches only itself), \texttt{eq\_of\_joinable\_normalForms}, and \texttt{uniqueNormalFormFrom\_of\_confluent}. The two directions of \Cref{thm:elementary-flatness-iff-local-confluence} are \texttt{elementaryFlat\_of\_localConfluent} and \texttt{localConfluent\_of\_elementaryFlat}.

\subsection{Axioms}

There are no \texttt{sorry} placeholders and no custom axioms. According to Lean's \texttt{\#print axioms}, the formal versions of \Cref{prop:confluence-normalization-unique-normal-form}, \Cref{thm:newman}, and \Cref{thm:normalizing-confluence-iff-unf}, as well as the supporting lemma that confluence gives unique reachable normal forms, depend on no axioms at all. The formal versions of \Cref{thm:elementary-flatness-iff-local-confluence}, \Cref{lem:terminating-normalizing}, \Cref{cor:terminating-confluence-iff-flat}, and \Cref{cor:terminating-flat-unique-nf} depend only on Lean's standard axioms \texttt{propext}, \texttt{Classical.choice}, and \texttt{Quot.sound}, which enter through classical case distinctions, such as whether two states are equal or whether a state has a successor.

\section{Computational Protocol}
\label{app:computational-protocol}

This appendix summarizes the computational protocol used to generate the evidence layers cited throughout the paper. The purpose of the protocol is reproducibility and claim discipline rather than algorithmic novelty.
The implementation repository for the experiments and code is \url{https://github.com/ioannist/six-birds-confluence}.

\subsection{Finite-ARS exhaustive protocol}

The finite-ARS audit enumerates every labeled directed rewriting relation on $n$ states for $n \leq 4$, including self-loops. This yields a total search space of
\[
2^1 + 2^4 + 2^9 + 2^{16} = 66066
\]
systems. Terminating systems are identified by directed-acyclicity of the finite reduction graph. For each terminating system, the repository computes local peak counts, nonjoinability defect counts, normal-form mismatch counts, local confluence, confluence, and uniqueness of reachable normal forms. The resulting summary and row-level artifacts are stored under \texttt{results/audits/}.

\subsection{Bounded string-rewriting protocol}

The string-rewriting examples are explored from designated start words under explicit depth, state-count, and word-length bounds. The exploration records a completeness flag so that the paper can distinguish complete explored closures from genuinely truncated ones. In the curated examples used here, the chosen bounds are sufficient for complete exploration. The repository computes graph-level local peaks, system critical pairs, reachable critical pairs relative to the designated starts, and the corresponding defect counts and signature comparisons.

\subsection{Bounded term-rewriting protocol}

The term-rewriting bridge uses a bounded first-order rewriting prototype with explicit start terms and exploration bounds. The critical-pair generator supports the left-linear overlap regime needed for the curated examples, including nonroot nonvariable-position overlaps. As in the string layer, the repository compares reachable critical-pair signatures with graph-peak signatures in the explored reduction graph and records whether the comparison is exact.

\subsection{Deterministic witness and filler selection}

Across the plaquette, critical-pair, and two-complex layers, witnesses and filler paths are selected deterministically. The policy is: minimize total branch-to-witness distance, break ties lexicographically on the witness, and then choose lexicographically smallest shortest paths from each branch to the chosen witness. This policy ensures that the machine-readable artifacts are stable across runs and that the resulting local generators can be compared reliably across domains.

\subsection{Regression and artifact generation}

The repository includes a regression harness that runs the unit tests, the curated finite-ARS checks, the curated structured rewriting checks, the small exhaustive audit, the diagram-generation step, and the results-ledger refresh. The frozen evidence boundary used in this paper is recorded by the final experiment summary, the claim-support matrix, and the central results ledger under \code{results/index.json} and \code{results/summary.csv}.

\section{Artifact Guide}
\label{app:artifact-guide}

Every number in the tables of this paper is read from a machine-readable record in the repository. This appendix lists where those records live.

\subsection{Records by family}

\begin{itemize}
    \item Finite abstract rewriting systems: the curated examples are in \code{data/examples/}; the exhaustive enumeration, the counterexample search, and the cycle-proxy comparison are in \code{results/audits/}; peak and filler records are in \code{results/peaks/} and \code{results/plaquettes/}.
    \item String and term rewriting: explored systems are in \code{results/string_rewrite/} and \code{results/term_rewrite/}, and critical-pair analyses, including the lists of reachable critical-pair instances, are in \code{results/critical_pairs/} and \code{results/term_rewrite/}.
    \item Reduction 2-complexes: records of $0$-cells, $1$-cells, and generating branchings, with filled or open status, join witnesses, and closed boundary walks, are in \code{results/two_complex/}; Graphviz diagram sources are in \code{results/diagrams/}.
    \item Completion pairs: the string pair is recorded as a case study in \code{results/case_studies/}; the term pair is recorded through its two explored systems in \code{results/term_rewrite/} and its branchings in \code{results/two_complex/}.
    \item Summary records: the claim-status matrix and the final experiment summary are in \code{results/freeze/}; \code{results/index.json} and \code{results/summary.csv} index all records and are the recommended starting point.
    \item Verification: \code{review/} contains the independent Warshall check, the Lean axiom report, and the record of the last full regression run.
\end{itemize}

\subsection{Sections and records}

\Cref{sec:bare-graph-not-enough} uses the cycle-proxy comparison and the counterexample search. \Cref{sec:reduction-two-complex} uses the reduction 2-complex records. \Cref{sec:critical-pairs-curvature} and \Cref{subsec:curated-string-systems,subsec:curated-trs-bridge} use the string, term, and critical-pair records. \Cref{subsec:exhaustive-finite-ars-audit,subsec:boundary-search} use the enumeration and counterexample records and the independent check. \Cref{sec:completion-holonomy} uses the completion and reduction 2-complex records. Earlier planning maps from sections and from table and figure slots to records are kept in \code{paper/notes/}; they predate some of the present tables and figures. The figures of this paper are drawn directly in its LaTeX sources from the examples they depict.

\section{Version History}
\label{app:version-history}

\paragraph{Version 1 (17 March 2026).} First release, DOI \href{https://doi.org/10.5281/zenodo.19061345}{10.5281/zenodo.19061345}.

\paragraph{Version 2 (1 October 2026).} DOI \href{https://doi.org/10.5281/zenodo.23088991}{10.5281/zenodo.23088991}.
\begin{itemize}
    \item The Lean development now covers arbitrary rewriting relations without finiteness assumptions, and adds termination, Newman's lemma, the equivalence of confluence and elementary flatness under termination, and the equivalence of confluence and unique normal forms under normalization (\Cref{sec:core-theorem}).
    \item The negative result on graph flatness is stated for the two cycle-based proxies; the reduction graph of a finite system does determine its confluence (\Cref{sec:bare-graph-not-enough}).
    \item Elementary holonomy is defined by joinability alone, and $2$-cells are attached only to filled branchings (\Cref{sec:reduction-two-complex}).
    \item Critical pairs are presented as generators of the nontrivial branchings rather than of all peaks, and the computations use concrete reachable critical-pair instances, including instances in context and under substitution (\Cref{sec:critical-pairs-curvature}).
    \item The computational pipeline was corrected: self-overlaps of term rules are included, unification performs a full occurs check, substitution is simultaneous, invalid rules are rejected, and the critical-pair and holonomy analyses refuse to certify verdicts from incomplete explorations (\Cref{app:computational-protocol}). All curated counts and verdicts are unchanged.
    \item New and redrawn figures, a table of smallest counterexamples, and a rewritten exposition.
\end{itemize}


\bibliographystyle{plain}
\bibliography{references}

\end{document}
